% claim.tex -- notchgrid-parity-mod priority claim note (SOP 08b claim lane).
% AI-generated 2026-09-28 by the AI4Math paper department (claim lane).
% Instantiated from harness/templates/claim/claim.tex; lstlean.tex copied unchanged
% from harness/templates/paper/lstlean.tex (single source of truth is the paper side).
% Machine-readable anchors: every theorem carries a "% LEAN:" line (source task id +
% Lean declaration name + grade); scripts/paper-lint.sh and the claim audit check
% against it. Grades are quoted from the final audit of the source task, not decided
% here: all four are 完全证明 (complete proof).
% All Lean listings are byte-exact mechanical sed excerpts of the final artifact
% tasks/20260927-mossad-notch/final/01-notch-proved.lean
% (sha256 6f1a7440f5830633e9309ab5629e16308928bb6421f64f4d4b1988b8929db4aa).
% The two definitions and the four theorem headers are character-identical to the
% frozen statement file tasks/20260927-mossad-notch/formalized/01-notch-mods.lean
% (sha256 9a787b6be05dc345343edf675d107801e674e3abae095356ace390fac8426fb8); verified
% independently in audit/statement-vs-frozen.txt. The frozen file carries "sorry"
% proof bodies, so no listing quotes anything beyond a declaration header.
% Wording ceiling (constitution clause 4 + C9 adjudication): the four theorems are
% about recurrence-defined integer sequences ONLY; the identification with matching
% counts of notched grids is computational evidence and is not formalized. The
% derived perfect-matching sub-problem for the 3-by-n notch is already OEIS A001353
% and is fenced off from any novelty statement.
\documentclass[11pt]{article}

\usepackage[T1]{fontenc}
\usepackage[utf8]{inputenc}
\usepackage{amsmath,amssymb,amsthm}
\usepackage{array}
\usepackage{hyperref}
% ragged-right table columns: the C9 channel table holds long unbreakable
% \texttt{...} query strings, which otherwise produce stretched lines.
\newcolumntype{L}[1]{>{\raggedright\arraybackslash}p{#1}}
% lstlean.tex — Lean style for the listings package.
% Lineage: Jeremy Avigad (2015), adapted from Assia Mahboubi's SSR style;
% inspired by Olivier Verdier's unixode.sty for the Unicode literate table.
% No CTAN package or upstream repo exists; papers carry this file in their
% source package (cap set arXiv:1907.01449, sphere eversion arXiv:2210.07746).
% AI4Math adaptation (AI-generated, 2026-09-26): sorry/admit/sorryAx elevated
% to a dedicated error tier, matching the pipeline's 0-sorry constitution.
\usepackage{listings}
\usepackage{xcolor}

\definecolor{keywordcolor}{rgb}{0.7, 0.1, 0.1}   % red keywords
\definecolor{tacticcolor}{rgb}{0.0, 0.1, 0.6}    % blue tactics
\definecolor{commentcolor}{rgb}{0.4, 0.4, 0.4}   % grey comments
\definecolor{symbolcolor}{rgb}{0.0, 0.1, 0.6}    % blue symbols
\definecolor{sortcolor}{rgb}{0.1, 0.5, 0.1}      % green sorts
\definecolor{errorcolor}{rgb}{0.6, 0, 0}         % dark red: sorry/admit
\definecolor{stringcolor}{rgb}{0.5, 0.1, 0.5}    % purple strings

\lstdefinelanguage{lean}{
  % Anything between $ becomes LaTeX math mode
  mathescape=false,
  % Comments may or not include Latex commands
  texcl=false,
  % keywords, list taken from lean-syntax.txt
  morekeywords=[1]{import, prelude, protected, private, noncomputable,
    definition, meta, renaming, hiding, parameter, parameters, begin,
    constant, constants, lemma, variable, variables, theory, print,
    theorem, example, open, section, namespace, universe, universes, end,
    modifier, modifiers, using, export, exposing, axiom, inductive,
    coinductive, with, structure, class, instance, attribute, local, where,
    by, have, show, let, in, if, then, else, match, do, for, fun, assume,
    from, take, calc, include, omit, infix, infixl, infixr, notation,
    macro, syntax, elab, set_option, initialize, deriving, opaque, abbrev,
    mutual, partial, scoped, termination_by, decreasing_by},
  % Sorts
  morekeywords=[2]{Sort, Type, Prop},
  % tactics, list taken from lean-syntax.txt (tactic keywords)
  morekeywords=[3]{intro, intros, apply, exact, refine, rfl, simp, simp_all,
    simpa, simp_rw, rw, rewrite, erewrite, ring, ring_nf, omega, decide,
    norm_num, norm_cast, induction, cases, rcases, obtain, constructor,
    ext, funext, conv, congr, field_simp, linarith, nlinarith, positivity,
    tauto, trivial, aesop, use, by_cases, by_contra, push_neg, contrapose,
    symm, trans, assumption, contradiction, exfalso, specialize,
    generalize, revert, clear, rename_i, subst, unfold, delta, change,
    convert, split_ifs, interval_cases, fin_cases, zify, gcongr,
    nth_rewrite, suffices, generalizing, at, only, native_decide},
  % warnings - constitution tier: must never survive into a final proof
  morekeywords=[4]{sorry, admit, sorryAx},
  literate=
    {α}{{\ensuremath{\alpha}}}1
    {β}{{\ensuremath{\beta}}}1
    {γ}{{\ensuremath{\gamma}}}1
    {δ}{{\ensuremath{\delta}}}1
    {ε}{{\ensuremath{\varepsilon}}}1
    {ζ}{{\ensuremath{\zeta}}}1
    {η}{{\ensuremath{\eta}}}1
    {θ}{{\ensuremath{\theta}}}1
    {ι}{{\ensuremath{\iota}}}1
    {κ}{{\ensuremath{\kappa}}}1
    {λ}{{\ensuremath{\lambda}}}1
    {μ}{{\ensuremath{\mu}}}1
    {ν}{{\ensuremath{\nu}}}1
    {ξ}{{\ensuremath{\xi}}}1
    {π}{{\ensuremath{\pi}}}1
    {ρ}{{\ensuremath{\rho}}}1
    {σ}{{\ensuremath{\sigma}}}1
    {τ}{{\ensuremath{\tau}}}1
    {υ}{{\ensuremath{\upsilon}}}1
    {φ}{{\ensuremath{\varphi}}}1
    {χ}{{\ensuremath{\chi}}}1
    {ψ}{{\ensuremath{\psi}}}1
    {ω}{{\ensuremath{\omega}}}1
    {Γ}{{\ensuremath{\Gamma}}}1
    {Δ}{{\ensuremath{\Delta}}}1
    {Θ}{{\ensuremath{\Theta}}}1
    {Λ}{{\ensuremath{\Lambda}}}1
    {Ξ}{{\ensuremath{\Xi}}}1
    {Π}{{\ensuremath{\Pi}}}1
    {Σ}{{\ensuremath{\Sigma}}}1
    {Φ}{{\ensuremath{\Phi}}}1
    {Ψ}{{\ensuremath{\Psi}}}1
    {Ω}{{\ensuremath{\Omega}}}1
    {ℵ}{{\ensuremath{\aleph}}}1
    {∀}{{\ensuremath{\forall}}}1
    {∃!}{{\ensuremath{\exists}!}}1
    {∃}{{\ensuremath{\exists}}}1
    {∈}{{\ensuremath{\in}}}1
    {∉}{{\ensuremath{\notin}}}1
    {∘}{{\ensuremath{\circ}}}1
    {∩}{{\ensuremath{\cap}}}1
    {∪}{{\ensuremath{\cup}}}1
    {⊆}{{\ensuremath{\subseteq}}}1
    {⊂}{{\ensuremath{\subset}}}1
    {⊇}{{\ensuremath{\supseteq}}}1
    {⊃}{{\ensuremath{\supset}}}1
    {⊄}{{\ensuremath{\nsubseteq}}}1
    {∅}{{\ensuremath{\emptyset}}}1
    {∧}{{\ensuremath{\wedge}}}1
    {∨}{{\ensuremath{\vee}}}1
    {¬}{{\ensuremath{\neg}}}1
    {⊤}{{\ensuremath{\top}}}1
    {⊥}{{\ensuremath{\bot}}}1
    {↔}{{\ensuremath{\leftrightarrow}}}1
    {→}{{\ensuremath{\rightarrow}}}1
    {←}{{\ensuremath{\leftarrow}}}1
    {↑}{{\ensuremath{\uparrow}}}1
    {⇒}{{\ensuremath{\Rightarrow}}}1
    {⇐}{{\ensuremath{\Leftarrow}}}1
    {⟹}{{\ensuremath{\Longrightarrow}}}1
    {↦}{{\ensuremath{\mapsto}}}1
    {≤}{{\ensuremath{\leq}}}1
    {≥}{{\ensuremath{\geq}}}1
    {≠}{{\ensuremath{\neq}}}1
    {≈}{{\ensuremath{\approx}}}1
    {≡}{{\ensuremath{\equiv}}}1
    {≃}{{\ensuremath{\simeq}}}1
    {≅}{{\ensuremath{\cong}}}1
    {×}{{\ensuremath{\times}}}1
    {·}{{\ensuremath{\cdot}}}1
    {±}{{\ensuremath{\pm}}}1
    {⊕}{{\ensuremath{\oplus}}}1
    {⊗}{{\ensuremath{\otimes}}}1
    {⋆}{{\ensuremath{\star}}}1
    {▸}{{\ensuremath{\blacktriangleright}}}1
    {⟨}{{\ensuremath{\langle}}}1
    {⟩}{{\ensuremath{\rangle}}}1
    {⦃}{{\ensuremath{\{\!\mid}}}1
    {⦄}{{\ensuremath{\mid\!\}}}}1
    {⟦}{{\ensuremath{[\![}}}1
    {⟧}{{\ensuremath{]\!]}}}1
    {⌊}{{\ensuremath{\lfloor}}}1
    {⌋}{{\ensuremath{\rfloor}}}1
    {⌈}{{\ensuremath{\lceil}}}1
    {⌉}{{\ensuremath{\rceil}}}1
    {√}{{\ensuremath{\surd}}}1
    {∣}{{\ensuremath{\mid}}}1
    {∤}{{\ensuremath{\nmid}}}1
    {∎}{{\ensuremath{\blacksquare}}}1
    {⋯}{{\ensuremath{\cdots}}}1
    {…}{{\ensuremath{\ldots}}}1
    {∑}{{\ensuremath{\sum}}}1
    {∏}{{\ensuremath{\prod}}}1
    {⁻¹}{{\ensuremath{^{-1}}}}1
    {¹}{{\ensuremath{^1}}}1
    {²}{{\ensuremath{^2}}}1
    {³}{{\ensuremath{^3}}}1
    {ⁿ}{{\ensuremath{^n}}}1
    {ᵐ}{{\ensuremath{^m}}}1
    {₀}{{\ensuremath{_0}}}1
    {₁}{{\ensuremath{_1}}}1
    {₂}{{\ensuremath{_2}}}1
    {₃}{{\ensuremath{_3}}}1
    {₄}{{\ensuremath{_4}}}1
    {₅}{{\ensuremath{_5}}}1
    {ₙ}{{\ensuremath{_n}}}1
    {ₘ}{{\ensuremath{_m}}}1
    {ₖ}{{\ensuremath{_k}}}1
    {ᵢ}{{\ensuremath{_i}}}1
    {ⱼ}{{\ensuremath{_j}}}1
    {ℤ}{{\ensuremath{\mathbb{Z}}}}1
    {ℕ}{{\ensuremath{\mathbb{N}}}}1
    {ℚ}{{\ensuremath{\mathbb{Q}}}}1
    {ℝ}{{\ensuremath{\mathbb{R}}}}1
    {ℂ}{{\ensuremath{\mathbb{C}}}}1
    {𝔽}{{\ensuremath{\mathbb{F}}}}1
    {𝒫}{{\ensuremath{\mathcal{P}}}}1
    {∗}{{\ensuremath{\ast}}}1
    {•}{{\ensuremath{\bullet}}}1,
  % Comments
  morecomment=[l]{--},
  morecomment=[s]{/-}{-/},
  morestring=[b]",
  stringstyle=\color{stringcolor},
  keepspaces=true,
  keywordstyle=[1]\color{keywordcolor}\bfseries,
  keywordstyle=[2]\color{sortcolor}\bfseries,
  keywordstyle=[3]\color{tacticcolor}\bfseries,
  keywordstyle=[4]\color{errorcolor}\bfseries\underline,
  escapeinside={(*@}{@*)},
  columns=fullflexible,
  basicstyle=\ttfamily\small,
  commentstyle=\color{commentcolor}\itshape,
  breaklines=true,
  showstringspaces=false,
  frame=single,
  framesep=2mm,
  xleftmargin=4mm,
  numbers=left,
  numberstyle=\tiny\color{commentcolor},
  numbersep=6pt,
}

% Inline Lean code in prose: \lean{Int.ModEq.pow}
\newcommand{\lean}[1]{\lstinline[language=lean,basicstyle=\ttfamily\normalsize]|#1|}

\lstset{language=lean}
% typography: let inter-word glue stretch rather than overflow lines holding
% unbreakable artifact path names.
\setlength{\emergencystretch}{2.5em}

% --- XeTeX 渲染补丁（模板自带，血统 papers/cyl3-mod23；非经用户同意不得删）---
% 作用：tectonic/XeTeX 下 listings literate 对多字节字符不生效（cyl3/edgemid 判例），
% 把 Lean 源码里的 Unicode 符号映射为数学字形（lccode 活动字符法，ASCII 源、不加宏包）；
% xeCJK+SimSun 渲染 listing 中文注释（本机实测 0 Missing character）。
% 仅影响渲染，listing 源字节不变（逐字节核验不受影响）；\ifdefined 守卫使 pdflatex（arXiv）路径零变化。
\ifdefined\XeTeXversion
  {\lccode`\~="2022 \lowercase{\global\catcode`~=\active \global\def~{\ensuremath{\bullet}}}}
  {\lccode`\~="2115 \lowercase{\global\catcode`~=\active \global\def~{\ensuremath{\mathbb{N}}}}}
  {\lccode`\~="2124 \lowercase{\global\catcode`~=\active \global\def~{\ensuremath{\mathbb{Z}}}}}
  {\lccode`\~="2190 \lowercase{\global\catcode`~=\active \global\def~{\ensuremath{\leftarrow}}}}
  {\lccode`\~="2192 \lowercase{\global\catcode`~=\active \global\def~{\ensuremath{\rightarrow}}}}
  {\lccode`\~="2194 \lowercase{\global\catcode`~=\active \global\def~{\ensuremath{\leftrightarrow}}}}
  {\lccode`\~="2200 \lowercase{\global\catcode`~=\active \global\def~{\ensuremath{\forall}}}}
  {\lccode`\~="2203 \lowercase{\global\catcode`~=\active \global\def~{\ensuremath{\exists}}}}
  {\lccode`\~="2227 \lowercase{\global\catcode`~=\active \global\def~{\ensuremath{\wedge}}}}
  {\lccode`\~="2228 \lowercase{\global\catcode`~=\active \global\def~{\ensuremath{\vee}}}}
  {\lccode`\~="2261 \lowercase{\global\catcode`~=\active \global\def~{\ensuremath{\equiv}}}}
  {\lccode`\~="2264 \lowercase{\global\catcode`~=\active \global\def~{\ensuremath{\leq}}}}
  {\lccode`\~="2265 \lowercase{\global\catcode`~=\active \global\def~{\ensuremath{\geq}}}}
  {\lccode`\~="22A2 \lowercase{\global\catcode`~=\active \global\def~{\ensuremath{\vdash}}}}
  {\lccode`\~="27E8 \lowercase{\global\catcode`~=\active \global\def~{\ensuremath{\langle}}}}
  {\lccode`\~="27E9 \lowercase{\global\catcode`~=\active \global\def~{\ensuremath{\rangle}}}}
\fi
\ifdefined\XeTeXversion
  \usepackage{xeCJK}
  \setCJKmainfont{SimSun}
\fi

\newtheorem{theorem}{Theorem}

\title{Parity and congruence periodicity of two integer recurrences obtained from
matching counts in notched $3 \times n$ and $4 \times n$ grids:\\
four theorems machine-checked in Lean 4 (claim of record)}
\author{Author Name\thanks{Contact and ORCID placeholder --- to be filled in by
the author before any upload.}}
\date{September 28, 2026}

\begin{document}
\maketitle

\begin{abstract}
We record four theorems, machine-checked in Lean 4, about two recurrence-defined
integer sequences \texttt{a3}, \texttt{a4} $\mathbb{N} \to \mathbb{Z}$, whose
recurrences --- order $5$ and order $9$ --- were obtained by exact enumeration of
the matchings of a $3 \times n$ and a $4 \times n$ grid graph with the top-right
vertex deleted. Namely: $\texttt{a3}(n)$ is odd exactly when $n \bmod 6 \in
\{2,3\}$; $\texttt{a4}(n)$ is odd exactly when $n \bmod 5 \in \{0,2\}$;
$\texttt{a3}(n+12) \equiv \texttt{a3}(n) \pmod 8$; and $\texttt{a4}(n) \equiv 2
\pmod 4$ exactly when $n \bmod 10 = 6$. All four proofs compile with no
placeholder and depend only on \{propext, Quot.sound, Classical.choice\}; the
statements are byte-identical to a hash-frozen file. A structured novelty search
found no occupying record for either sequence (adjudicated value tier: new
sequence / new recurrence); the perfect-matching sub-problem of the $3 \times n$
notch is excluded from that because it is already OEIS A001353. The
identification of the sequences with matching counts is \emph{not} formalized.
This note is a priority claim of record, not a full paper.
\end{abstract}

\section{Introduction}

Two graph families are considered: the grid graph $P_3 \square P_n$ and the grid
graph $P_4 \square P_n$, each with its top-right vertex deleted (the ``notch'').
Their matchings (edge sets with no two edges sharing an endpoint, the empty set
included) are counted two ways: a column-frontier dynamic program, which produced
the whole $24$-row count table, and a brute-force recursion over matchings (the
least uncovered vertex left unmatched or paired with one of its neighbours), run
over the leading eight and six entries respectively --- there the two methods agree
term by term. Berlekamp--Massey arithmetic over $\mathbb{Q}$ on those terms yields
two integer linear recurrences, of order $5$ for the $3 \times n$ family and
order $9$ for the $4 \times n$ family. Taking those recurrences and their initial
values as \emph{definitions}, this note records four theorems about the parities
and residue classes of the two sequences, checked by the Lean 4 kernel.

The work was produced by AI4Math, a budgeted LLM-plus-kernel pipeline in which
compilation is the sole adjudicator at every gate. Each theorem below carries a
machine-readable anchor comment tying it to its Lean declaration name and to the
auditor's grade, \emph{complete proof} for all four --- quoted from the audit
record, not a judgement made in this note.

This text is a priority claim of record: a short deposit note accompanying a
Zenodo package, so that the two sequences and their recurrences are publicly
timestamped while the slower paper lane proceeds (a fuller narrative paper on the
same result, at \texttt{papers/notchgrid-parity-mod/}, will cite this deposit).
Repository paths quoted in this note are internal ones: they are not publicly
retrievable at deposit time, and the public mirror of this record is exactly the
contents of this package. A literature survey, a prose proof and a research-level
narrative are not duties of this note and are not attempted here.

\section{Statement}\label{sec:stmt}

Both sequences are functions $\mathbb{N} \to \mathbb{Z}$, not $\mathbb{N} \to
\mathbb{N}$: the recurrences have negative coefficients, so the integer codomain
lets the defining equations be read as ordinary subtraction on $\mathbb{Z}$, and
it lets the \texttt{Int.ModEq} congruences and the \texttt{Odd} bridge of
Section~\ref{sec:proof} apply without further coercion. Both are $0$-indexed, so
that $\texttt{a3}(0) = 2$ and $\texttt{a4}(0) = 3$ are the notch on one column.
The two definitions, quoted verbatim from the artifact and character-identical to
the frozen statement file, are:

\begin{lstlisting}[caption={\texttt{a3}: five initial values, order-5 recurrence
(verbatim excerpt, lines 24--31 of \texttt{final/01-notch-proved.lean}).},
label={lst:a3def}]
def a3 : ℕ → ℤ
  | 0 => 2
  | 1 => 10
  | 2 => 67
  | 3 => 407
  | 4 => 2546
  | n + 5 =>
      5 * a3 (n + 4) + 9 * a3 (n + 3) - 9 * a3 (n + 2) - a3 (n + 1) + a3 n
\end{lstlisting}

\begin{lstlisting}[caption={\texttt{a4}: nine initial values, order-9 recurrence
(verbatim excerpt, lines 151--163). The order 9 is reported as found: it exceeds
the order 3--6 the external proposal anticipated.},label={lst:a4def}]
def a4 : ℕ → ℤ
  | 0 => 3
  | 1 => 32
  | 2 => 407
  | 3 => 4840
  | 4 => 58608
  | 5 => 705949
  | 6 => 8515850
  | 7 => 102684287
  | 8 => 1238310540
  | n + 9 =>
      9 * a4 (n + 8) + 41 * a4 (n + 7) - 41 * a4 (n + 6) - 111 * a4 (n + 5)
        + 91 * a4 (n + 4) + 29 * a4 (n + 3) - 23 * a4 (n + 2) - a4 (n + 1) + a4 n
\end{lstlisting}

Both recurrences were fitted on 21 terms with 3 held out (residuals zero), refit
on the minimal $2L$-term window giving the same recurrence, and re-substituted
against the independent dynamic program for all $n \leq 60$: computed evidence
throughout. The theorems take the recurrences as definitions and are universal
over $\mathbb{N}$.

% LEAN: tasks/20260927-mossad-notch :: a3_odd_iff grade=完全证明
\begin{theorem}[parity of \texttt{a3}]\label{thm:t1}
For every $n \in \mathbb{N}$, $\texttt{a3}(n)$ is odd if and only if $n \bmod 6 =
2$ or $n \bmod 6 = 3$.
\end{theorem}

\begin{lstlisting}[caption={Formal statement of Theorem~\ref{thm:t1}: declaration
header (verbatim excerpt, line 34).},label={lst:t1}]
theorem a3_odd_iff (n : ℕ) : Odd (a3 n) ↔ n % 6 = 2 ∨ n % 6 = 3 := by
\end{lstlisting}

% LEAN: tasks/20260927-mossad-notch :: a4_odd_iff grade=完全证明
\begin{theorem}[parity of \texttt{a4}]\label{thm:t2}
For every $n \in \mathbb{N}$, $\texttt{a4}(n)$ is odd if and only if $n \bmod 5 =
0$ or $n \bmod 5 = 2$.
\end{theorem}

\begin{lstlisting}[caption={Formal statement of Theorem~\ref{thm:t2}: declaration
header (verbatim excerpt, line 166).},label={lst:t2}]
theorem a4_odd_iff (n : ℕ) : Odd (a4 n) ↔ n % 5 = 0 ∨ n % 5 = 2 := by
\end{lstlisting}

% LEAN: tasks/20260927-mossad-notch :: a3_mod8_periodic grade=完全证明
\begin{theorem}[\texttt{a3} modulo 8 repeats with period 12]\label{thm:t3}
For every $n \in \mathbb{N}$, $\texttt{a3}(n+12) \equiv \texttt{a3}(n) \pmod 8$.
\end{theorem}

\begin{lstlisting}[caption={Formal statement of Theorem~\ref{thm:t3}: declaration
header (verbatim excerpt, line 122).},label={lst:t3}]
theorem a3_mod8_periodic (n : ℕ) : a3 (n + 12) ≡ a3 n [ZMOD 8] := by
\end{lstlisting}

% LEAN: tasks/20260927-mossad-notch :: a4_mod4_eq2_iff grade=完全证明
\begin{theorem}[\texttt{a4} is 2 modulo 4 in exactly one residue class]\label{thm:t4}
For every $n \in \mathbb{N}$, $\texttt{a4}(n) \equiv 2 \pmod 4$ if and only if
$n \bmod 10 = 6$.
\end{theorem}

\begin{lstlisting}[caption={Formal statement of Theorem~\ref{thm:t4}: declaration
header (verbatim excerpt, line 281).},label={lst:t4}]
theorem a4_mod4_eq2_iff (n : ℕ) : a4 n ≡ 2 [ZMOD 4] ↔ n % 10 = 6 := by
\end{lstlisting}

The residue tables over one full period starting at $n = 0$, computed from the
definitions by exact integer iteration ($n \leq 400$), are

\begin{center}\footnotesize
\setlength{\tabcolsep}{4pt}
\begin{tabular}{ll@{\hspace{2.5em}}ll}
$\texttt{a3} \bmod 2$ & $(0,0,1,1,0,0)$ &
$\texttt{a3} \bmod 8$ & $(2,2,3,7,2,6,0,0,7,7,0,0)$ \\
$\texttt{a4} \bmod 2$ & $(1,0,1,0,0)$ &
$\texttt{a4} \bmod 4$ & $(3,0,3,0,0,1,2,3,0,0)$
\end{tabular}
\end{center}

\noindent and they are also the pattern the Lean proofs work with: in
Theorems~\ref{thm:t1}, \ref{thm:t2} and \ref{thm:t4} the table appears as an
inner lemma, whereas Theorem~\ref{thm:t3} is proved as a shift identity and
contains no such table. Minimality of the periods $6$, $5$, $12$, $10$ is a
computed fact and is \emph{not} among the formalized statements;
Theorem~\ref{thm:t3} asserts periodicity \emph{with} period 12.

\section{Proof}\label{sec:proof}

The complete development is \texttt{final/}\allowbreak\texttt{01-notch-proved.%
lean} (621 lines), shipped in full in this package at
\texttt{proofs/01-notch-proved.lean}; the listing below is an \emph{excerpt} of
it. Three of the four proofs (Theorems~\ref{thm:t1}, \ref{thm:t2} and
\ref{thm:t4}) share one residue-table skeleton: establish an inner statement (a
\lean{have} step, never an extra top-level theorem) that pins the residue of
every term through an explicit \texttt{if}-chain over $n \bmod P$; enumerate the
base window below the recurrence order with \lean{interval_cases} and
\lean{decide}; in the inductive step write $n = j + L$ and unfold the recurrence
exactly once by \lean{simp only [a3]} or \lean{simp only [a4]}; split on
\lean{mod_cases} $j \bmod P$, so that in each branch the table value of every
predecessor is fixed and the weighted sum of the resulting congruences closes by
one \lean{decide}; then, in the two parity theorems, bridge back to \lean{Odd}
through \lean{Int.odd_iff}. Listing~\ref{lst:t3proof} gives the proof of
Theorem~\ref{thm:t3} complete, the most economical of the four, and it is
\emph{not} of that skeleton: no residue table, no \texttt{if}-chain, no
\lean{mod_cases} split, no \lean{Int.odd_iff} bridge. It is a direct congruence
calculation --- one \lean{simp only [a3]} step gives the recurrence equation,
used at two offsets to expand both $\texttt{a3}(j+12+12)$ and
$\texttt{a3}(j+12)$; the induction hypothesis applied to the five predecessors
one by one; the five congruences combined through the compatibility of
\texttt{Int.ModEq} with addition, subtraction and integer multiplication;
closing by rewriting with the two expansions, the base window being $n<12$
rather than the recurrence order. For Theorem~\ref{thm:t1} the artifact
shows the residue-table skeleton at its most explicit (inner residue-table lemma
with its induction
setup at lines 35--45, one of the six residue branches at lines 51--60, bridge
back to \texttt{Odd} at lines 106--119). Theorem~\ref{thm:t2} is
Theorem~\ref{thm:t1} with nine predecessors and period 5 (its body is lines
167--277); Theorem~\ref{thm:t4} is the case where the table is genuinely
four-valued, written as a three-level \texttt{if}-chain (inner lemma from line
282, ten forward branches to line 618, backward direction lines 619--621). Every
numeric evaluation inside these proofs, the base-window values of
Theorem~\ref{thm:t3} among them, is performed by the kernel's own \lean{decide}
on closed integer expressions; there is no external oracle and no
verification bound.

\begin{lstlisting}[caption={Proof of Theorem~\ref{thm:t3} in full (verbatim
excerpt, lines 123--144; the complete source of all four proofs is in
\texttt{proofs/01-notch-proved.lean}).},label={lst:t3proof}]
  induction n using Nat.strong_induction_on with
  | _ n ih =>
    rcases lt_or_ge n 12 with hlt | hge
    · interval_cases n <;> decide
    · obtain ⟨j, rfl⟩ : ∃ j, n = j + 12 := ⟨n - 12, by omega⟩
      have rec5 : ∀ m : ℕ, a3 (m + 5)
          = 5 * a3 (m + 4) + 9 * a3 (m + 3) - 9 * a3 (m + 2) - a3 (m + 1) + a3 m := fun m => by
        simp only [a3]
      have e1 : a3 (j + 12 + 12)
          = 5 * a3 (j + 23) + 9 * a3 (j + 22) - 9 * a3 (j + 21) - a3 (j + 20) + a3 (j + 19) :=
        rec5 (j + 19)
      have e2 : a3 (j + 12)
          = 5 * a3 (j + 11) + 9 * a3 (j + 10) - 9 * a3 (j + 9) - a3 (j + 8) + a3 (j + 7) :=
        rec5 (j + 7)
      have h23 : a3 (j + 23) ≡ a3 (j + 11) [ZMOD 8] := ih (j + 11) (by omega)
      have h22 : a3 (j + 22) ≡ a3 (j + 10) [ZMOD 8] := ih (j + 10) (by omega)
      have h21 : a3 (j + 21) ≡ a3 (j + 9) [ZMOD 8] := ih (j + 9) (by omega)
      have h20 : a3 (j + 20) ≡ a3 (j + 8) [ZMOD 8] := ih (j + 8) (by omega)
      have h19 : a3 (j + 19) ≡ a3 (j + 7) [ZMOD 8] := ih (j + 7) (by omega)
      have key := ((((h23.mul_left 5).add (h22.mul_left 9)).sub (h21.mul_left 9)).sub h20).add h19
      rw [e1, e2]
      exact key
\end{lstlisting}

\section{Verification evidence}

\begin{itemize}
\item Toolchain: Lean v4.34.0 (\texttt{leanprover/lean4:v4.34.0}) \cite{lean4},
mathlib4 v4.34.0, revision \texttt{5ed29652} \cite{mathlib}; the pin is shipped
verbatim as \texttt{proofs/lean-toolchain}.
\item Statement integrity: the frozen statement file is
\texttt{formalized/}\allowbreak\texttt{01-notch-mods.lean} (69 lines) and the
delivered development is \texttt{final/}\allowbreak\texttt{01-notch-proved.lean}
(621 lines); their sha256 values, which every listing in this note is checked
against, are
\begin{center}\footnotesize\ttfamily
9a787b6be05dc345343edf675d107801e674e3abae095356ace390fac8426fb8 \\
6f1a7440f5830633e9309ab5629e16308928bb6421f64f4d4b1988b8929db4aa
\end{center}
The symmetric difference after stripping proof bodies is empty: the 6 top-level
declarations (2 definitions and 4 theorem headers, names and order) are
character-identical, and the delivered file adds no top-level statement.
\item Kernel check: strict compile exits 0 with zero errors (320 deprecation
warnings about \texttt{if\_neg} and \texttt{if\_pos}, recorded by the auditor as
non-adjudicating); a byte-level text scan including comments finds no
\texttt{sorry} and no \texttt{admit}; an escape-hatch scan (\texttt{set\_option},
\texttt{axiom}, \texttt{unsafe}, \texttt{implemented\_by}, \texttt{extern},
\texttt{macro\_rules}, \texttt{elab}, \texttt{native\_decide}, \texttt{decide!})
finds none. Axiom printout, copied verbatim and in the order the file prints it:

{\scriptsize
\begin{verbatim}
'a3_odd_iff' depends on axioms: [propext, Classical.choice, Quot.sound]
'a3_mod8_periodic' depends on axioms: [propext, Classical.choice, Quot.sound]
'a4_odd_iff' depends on axioms: [propext, Classical.choice, Quot.sound]
'a4_mod4_eq2_iff' depends on axioms: [propext, Classical.choice, Quot.sound]
\end{verbatim}
}
\item Audit chain: gate 2 well-definedness (four gate-batch compiles accept the
statements; $8 \times 4 = 32$ bare decision-procedure probes all fail, so no
theorem is closable by one automation step; both sequences re-derived from the
frozen definitions agree with the 24-row enumeration table in both columns and
the four statements hold pointwise for $n = 0, \dots, 80$) and gate 3 final audit
(symmetric difference, \texttt{sorry} and escape-hatch scans, re-run of the
strict compile and of the axiom printout, both outputs byte-identical to the
worker's own logs).
What \texttt{audit/} of this package ships is the key-line extract of each gate,
the verbatim axiom printout and the claim lane's own re-verification files; the
gate 2 and gate 3 originals (\texttt{final-01-notch.txt},
\texttt{welldef-verdict.md}, \texttt{gate-t1.txt}--\texttt{gate-t4.txt},
\texttt{review-final-\{strict,axioms\}.out}) stay in the source task directory
\texttt{tasks/20260927-mossad-notch/} and are not part of this deposit.
\end{itemize}

\section{Novelty evidence}\label{sec:novelty}

The pipeline's C9 novelty review adjudicated the two main propositions as
\emph{no occupying record found} (on its four-state scale: occupied / no
occupying record found / known but not formalized / undecided), with value tier
\emph{new sequence / new recurrence} (a translated pool-card tier label; the
Chinese original is in \texttt{audit/c9-record.md} of this package). This is a
statement about search evidence, not a claim to new mathematics, and not a
priority claim beyond the deposit itself. Per-channel summary, one line per
channel; the full log, including every zero-hit query, is copied verbatim into
\texttt{audit/c9-record.md}:

\begin{center}
\footnotesize
\setlength{\tabcolsep}{4pt}
\begin{tabular}{L{0.17\textwidth}L{0.52\textwidth}L{0.21\textwidth}}
\hline
channel & query & result \\
\hline
OEIS, sequence layer & \texttt{a3} head \texttt{2,10,67,407,2546,15782,98104,%
609304} and \texttt{a4} head \texttt{3,32,407,4840,58608,705949}, each with four
variants (drop 1, drop 2, $\times 2$, $+1$); divide-by-2 not applicable &
all 10 queries zero-hit \\
OEIS, keyword layer & \texttt{Hosoya index grid graph} (3 hits), \texttt{monomer
dimer grid} (10 hits) & no notched-count entry \\
OEIS, triangle check & A210662 rows and columns cut by hand: $T(3,\cdot) =
3,22,131$; $T(n,3) = 131,823,5096,\dots$; $T(4,\cdot) = 5,71,823,10012$
\cite{oeis-a210662} & our heads are not segments \\
OEIS, derived layer & \texttt{a3} near-perfect counts $4,29,161,798$ (even
widths); \texttt{a4} near-perfect $2,7,29,88,288$ and its drop-1 tail & zero-hit
\\
OEIS, derived layer (fence) & \texttt{a3} \emph{perfect}-matching counts, odd
widths: $1,4,15,56,209,780,2911$ & \textbf{A001353 occupied} \\
literature, zbMATH & 5 query families: \texttt{monomer dimer grid} (12),
\texttt{matching polynomial grid graph} (90), \texttt{matchings grid graph
corner} (1), \texttt{deficient graph matchings} (2), \texttt{monomer dimer
lattice vacancy} (1) & no same proposition \\
literature, other & Math StackExchange \texttt{number of matchings grid graph}
(6 items); arXiv \texttt{monomer dimer grid graph} (0) & none matching \\
literature, degraded & OpenAlex: HTTP 429 throughout (unverified, contributes no
evidence either way); general web search: no search tool in that session, one
engine DNS-poisoned, the other SEO noise & literature layer rests on zbMATH \\
library layer & mathlib (no \texttt{GridGraph}; \texttt{SimpleGraph/Matching.%
lean} has matching predicates but no counting theorem), compfiles, sequencelib &
zero hits \\
Lean ecosystem layer & GitHub repository search for Lean graph-matching
formalizations & zero hits \\
channel self-check & \texttt{0,1,1,2,3,5,8,13} gives A000045;
\texttt{1,3,13,22,44,90,196,406} gives A180970 & retrieval confirmed working \\
\hline
\end{tabular}
\end{center}

Two fences belong to the adjudication and are stated here, not in a footnote.
(i) The \emph{perfect}-matching sub-problem of the $3 \times n$ notch (odd
widths) is OEIS \textbf{A001353}, $a(n) = 4a(n-1) - a(n-2)$ with $a(0) = 0$,
$a(1) = 1$, whose record carries the independent attribution to Joshua Zucker and
the Castilleja School Math Club (2003-10-28) for packing a $3 \times (2n-1)$
rectangle with one extra square attached \cite{oeis-a001353}; nothing here claims
that sub-problem and no novelty statement extends to it. Our object is the
\emph{total} matching count. (ii) For the $4 \times n$ notch the vertex count
$4n - 1$ is odd, so its perfect-matching count is identically zero: a degeneracy
noted only to explain why that direction is empty, not a proposition.

\section{Scope and limitations}

\begin{enumerate}
\item \textbf{Grading boundary (quoted, not decided here).} The audit record
grades all four theorems \emph{complete proof}: no \texttt{sorry}, axioms inside
the standard whitelist, statement byte-identical to the frozen text. That grade
is the wording ceiling of this note.
\item \textbf{The semantic gap.} What is machine-checked is a property of two
recurrence-defined integer sequences. That they \emph{are} the matching counts of
notched grids is computed evidence: the brute-force recursion agrees with the
dynamic program term by term over the leading $8$ entries of the $3 \times n$
column and the leading $6$ of the $4 \times n$ column, which is the whole window it
was run, so the remaining rows of the $24$-row table are produced by the dynamic
program alone; recurrence re-substitution was checked up to $n = 60$ and modular
dynamic programming up to $n = 600$. The bridge is stated as a conjecture in the
accompanying paper, not proved here. No sentence of this note should be read as a
theorem of graph enumeration.
\item \textbf{Novelty is search evidence only}, at the adjudicated tier and with
the two fences of Section~\ref{sec:novelty}; OpenAlex returned HTTP 429 and the
general web layer was degraded, so if either later returns a strong hit the
novelty sentence must be re-examined.
\item \textbf{The order-9 recurrence is reported as found.} It exceeds the order
$3$--$6$ the external proposal anticipated; the refit on the minimal 18-term
window gives the same recurrence and it holds on the whole independent window, so
it is not an artifact of a short fitting window: nothing was trimmed to make it
look smaller.
\item \textbf{Computed versus checked.} Irreducibility of the two characteristic
polynomials over $\mathbb{Q}$, minimality of the four periods, and every
verification window are computed evidence (Python, exact integer or rational
arithmetic), not Lean theorems.
\item \textbf{Provenance.} Proof search consumed no LLM sampling: the four proofs
follow a skeleton already taken through the kernel on a sibling task of the same
selection pool, and the auditor ruled that this lineage does not downgrade the
grade. The autoformalization roundtrip check ran on the same model node as the
formalizer, so it is not an independent review; the mitigation is that the kernel
is the final arbiter and that the audit re-ran both compile and axiom printout.
\item This note contains no literature survey and no prose proof: it is a deposit
record.
\end{enumerate}

\section*{Data availability}

Artifact package: Zenodo deposit (DOI to be reserved at upload; see
\texttt{UPLOAD.md} in the source repository) containing this note (source and
PDF),
\texttt{proofs/} (frozen statement file, delivered development, toolchain pin),
\texttt{audit/} (gate and compile evidence key lines, axiom printout,
symmetric-difference conclusion, listing byte-verification reports, the claim
audit verdict \texttt{claim-check.md}, the full C9 query record
\texttt{c9-record.md}) and \texttt{metadata/}. A sanitized public mirror of
exactly this package may be pushed to GitHub by the author, with the URL recorded
in \texttt{claims/notchgrid-parity-mod/card.md} upon push; no other part of the
pipeline is mirrored. The numerical layer that produced the two sequences
(enumeration, fitting and modular probing scripts, 24-row count table) lives under
\texttt{phase0/} of the source task directory, is not part of this deposit, and is
likewise not publicly retrievable. Note under CC BY 4.0, code under MIT.

\section*{Use of generative AI}

(a) \emph{Stages done by AI}: candidate-proposition generation, the counting and
fitting scripts, autoformalization into Lean statements, proof search, audit
re-runs, prose translation and the drafting of this note were performed by LLM
stages of the AI4Math pipeline; topic adjudication, statement-semantics
verification and gate sign-off were human. (b) \emph{Model, node and budget}:
mathematical content was sampled on one resolved node, wire-id
\texttt{anthropic/}\allowbreak\texttt{a6api-main/}\allowbreak\texttt{kimi-k3}
(resolution level L3 static profile cross-checked against L2 relay state, taken
2026-09-27), with no cross-node switching. Numbers copied verbatim from the
source ledger \texttt{tasks/20260927-mossad-notch/}\allowbreak\texttt{budget.log}:
llm-call $2/8$ (both calls are autoformalization roundtrip
re-translation checks, tokens in/out $1628/4000$ truncated and $1628/1889$
complete) and Lean compile $17/24$ (gate 2: 6, proof search: 9, final audit
re-runs: 2); Phase 0, proof search, audit and reporting used no sampling at all.
This claim-note stage used llm-call $0$ of its budget of $4$ and $9$ LaTeX
compile rounds of its budget of $9$ (the original packaging used the then cap of
$6$, which the user raised to $9$ on 2026-09-28 for the rework done after the
audit; the ninth round was spent on a package-promise wording correction, and
the DOI backfill, which needs one further round, is to be requested separately
under explicit user authorization; no cap was exceeded), as recorded in
\texttt{claims/notchgrid-parity-mod/budget.log}. The drafting session ran on a
different endpoint whose model selection key is \texttt{auto} (routing not
pinned); it contributes no sampled claim, and every assertion here that something
was ``proved'' or ``passed'' rests on kernel compilation or on the static gates,
not on any model's self-report. (c) \emph{Final arbiter}: all statements and
proofs reported here were checked by the Lean 4 kernel: every proof compiles with
no \texttt{sorry}, and \texttt{\#print axioms} reports only \{propext, Quot.sound,
Classical.choice\}. (d) The human author reviewed and is responsible for all
content; the AI system is not an author.

\section*{Author contributions}

CRediT-style: the human author is responsible for topic adjudication and
selection, verification of statement semantics against the frozen files, gate
sign-off, and the decision to deposit; the AI4Math pipeline (an LLM-based system)
carried out the enumeration scripts, autoformalization, proof search, audit
re-runs and drafting of this note. The AI system is not listed as an author.

\begin{thebibliography}{9}
% Every entry below was verified reachable online from this machine with title and
% authors cross-checked against doi.org / the Crossref API / the OEIS page itself:
% first at paper time (papers/notchgrid-parity-mod/audit/references-check.txt,
% 2026-09-27) and re-run at claim time
% (claims/notchgrid-parity-mod/audit/references-check.txt, 2026-09-28). Entries are
% copied verbatim from that verified bibliography; no entry was written from
% memory. The one paper entry not cited in this note (heilmann-lieb) was
% deliberately not carried over.
\bibitem{lean4}
L.~de Moura and S.~Ullrich.
\newblock The Lean 4 theorem prover and programming language.
\newblock In \emph{Automated Deduction --- CADE-28}, LNCS~12699, pp.~625--635,
2021.
\newblock \url{https://doi.org/10.1007/978-3-030-79876-5_37}

\bibitem{mathlib}
{The mathlib Community}.
\newblock The Lean mathematical library.
\newblock In \emph{Proceedings of the 9th ACM SIGPLAN International Conference
on Certified Programs and Proofs (CPP '20)}, pp.~367--381, 2020.
\newblock \url{https://doi.org/10.1145/3372885.3373824}

\bibitem{oeis-a001353}
{The OEIS Foundation}.
\newblock Entry A001353: $a(n) = 4a(n-1) - a(n-2)$ with $a(0)=0$ and $a(1)=1$;
the entry's comment of Joshua Zucker and the Castilleja School Math Club
(2003-10-28) on packing a $3 \times (2n-1)$ rectangle with one extra square
attached.
\newblock \emph{The On-Line Encyclopedia of Integer Sequences}.
\newblock \url{https://oeis.org/A001353}

\bibitem{oeis-a210662}
{The OEIS Foundation}.
\newblock Entry A210662: triangle $T(n,k)$, number of monomer-dimer tilings of
an $n \times k$ board.
\newblock \emph{The On-Line Encyclopedia of Integer Sequences}.
\newblock \url{https://oeis.org/A210662}
\end{thebibliography}

\end{document}
